% !TEX root =./ECE448_Practice_main.tex
%
% L13 practice set (DTMF). Every question is a step in the L13 receiver:
% reading the keypad grid, the timing, setting a Goertzel block at another
% length, reading twist off the display, the decoder's decision rules, and
% what the link shows as noise goes up. No error-rate or bin math (Neil,
% 2026-10-21).
% Blank copy: \begin{solution}[h] reserves h of work space.
% SOLUTIONS copy (\iskey=1) prints the red worked solutions.

\fullwidth{\textbf{Learning Objective 6.2: } I can decode DTMF tones from a captured audio
signal.
\\[4pt]\normalfont\small Show your work. Write a \textbf{Documentation} line at the
end (\textbf{None} if you did not collaborate). The keypad grid: rows 697, 770, 852, and
941 Hz carry 1 2 3 A, 4 5 6 B, 7 8 9 C, and * 0 \# D; columns are 1209, 1336, 1477, and
1633 Hz, left to right.}

\begin{questions}

%%%%%%%%%%%%%%%%%%%%%%%%%%%%%%%%%%%%%%%%%
\question \textbf{Reading the grid.}
\begin{parts}

	\begin{minipage}[t]{0.58\linewidth}
		\part Which two tones does the `8' key send? The `\#' key?
		\begin{solution}[0.7in]
		`8' is row 3 and column 2: 852 Hz and 1336 Hz. `\#' is row 4 and column 3: 941 Hz
		and 1477 Hz.
		\end{solution}
	\end{minipage}\hfill%
	\begin{minipage}[t]{0.37\linewidth}
		\raggedleft \vspace{12pt}
		8: \ansbox[1.3in]{852, 1336 Hz}

		\vspace{4pt}
		\#: \ansbox[1.3in]{941, 1477 Hz}
	\end{minipage}

	\begin{minipage}[t]{0.58\linewidth}
		\part A receiver's bins show 941 Hz and 1209 Hz in one digit and 770 Hz and
		1633 Hz in the next. Which two keys were pressed?
		\begin{solution}[0.7in]
		941 and 1209 Hz is row 4, column 1: `*'. 770 and 1633 Hz is row 2, column 4: `B'.
		\end{solution}
	\end{minipage}\hfill%
	\begin{minipage}[t]{0.37\linewidth}
		\raggedleft \vspace{12pt}
		\ansbox[1.1in]{*, B}
	\end{minipage}

	\part In one block, the 697 Hz and 770 Hz bins are both strong and about equal, and so
	is the 1336 Hz bin. Which test of the L13 decoder rejects the block, and what could
	cause a signal like that?
		\begin{solution}[0.9in]
		Test 2: the best row bin must beat the other three row bins by 6 dB, and here two row
		bins are equal. A keypad never sends two row tones at once, so the signal is something
		else: two keys pressed together, or speech or music with energy near both
		frequencies.
		\end{solution}

\end{parts}

\newpage
%%%%%%%%%%%%%%%%%%%%%%%%%%%%%%%%%%%%%%%%%
\question \textbf{Timing.} An auto-dialer sends each digit as 70 ms of tone followed by
15 ms of silence.
\begin{parts}

	\begin{minipage}[t]{0.58\linewidth}
		\part How long does one digit take, tone plus gap, and how long does an 11-digit
		number take, from the first tone to the end of the last one?
		\begin{solution}[0.9in]
		Each digit takes $70 + 15 = 85$ ms.
		\eq{ 11 \times 85 - 15 = 920\ \text{ms} }
		\end{solution}
	\end{minipage}\hfill%
	\begin{minipage}[t]{0.37\linewidth}
		\raggedleft \vspace{12pt}
		digit: \ansbox[1.1in]{85 ms}

		\vspace{4pt}
		number: \ansbox[1.1in]{0.92 s}
	\end{minipage}

	\part Lesson 13 says DTMF tones last 40 ms or longer and gaps 20 ms or longer. Which
	does this dialer break? You set the same tone and gap on the sliders of
	\texttt{L13-dtmf-link.grc} and dial 7195550148. What would you expect to see in the
	Digits label?
		\begin{solution}[1.0in]
		The 70 ms tones are fine, but the 15 ms gaps are shorter than 20 ms. A gap that short
		can fall between the decoder's blocks, so a repeated digit merges: some dialings read
		71950148 or 719550148, with the three 5s counted as one or two.
		\end{solution}

\end{parts}

\newpage
%%%%%%%%%%%%%%%%%%%%%%%%%%%%%%%%%%%%%%%%%
\question \textbf{A shorter Goertzel block.} You want to try Goertzel blocks of $N = 136$
samples at $f_s = 8$ kHz in \texttt{L13-dtmf-rx.grc} instead of the lesson's 205.
\begin{parts}

	\begin{minipage}[t]{0.58\linewidth}
		\part Find the block time and the bin spacing, which is also the distance from a bin to its first null.
		\begin{solution}[0.8in]
		\eq{ N / f_s &= 136 / 8000 = 17\ \text{ms} \\
			f_s / N &= 8000 / 136 = 58.8\ \text{Hz} }
		\end{solution}
	\end{minipage}\hfill%
	\begin{minipage}[t]{0.37\linewidth}
		\raggedleft \vspace{12pt}
		block: \ansbox[1.1in]{17 ms}

		\vspace{4pt}
		spacing: \ansbox[1.1in]{58.8 Hz}
	\end{minipage}

	\part What do you change in the flowgraph, and what happens to the eight Goertzel
	blocks' Length and the two time sinks that plot the bins?
		\begin{solution}[0.8in]
		Set the Variable \texttt{N} to 136. The eight Goertzel lengths and the
		\texttt{samp\_rate/N} sample rate of the time sinks follow it, so nothing else
		changes; the sinks now get about 58.8 outputs a second.
		\end{solution}

	\part Compared with $N = 205$, what do you expect the shorter block to do better, and
	what worse? Answer in words.
		\begin{solution}[1.0in]
		Better: it responds faster, so it catches shorter tones, and its wider bins still
		decode a tone that is a little off frequency. Worse: its wider bins let more of the
		neighboring tone and more noise in, so the right bin wins by less, and short bursts
		of speech are more likely to look like a digit.
		\end{solution}

\end{parts}

\newpage
%%%%%%%%%%%%%%%%%%%%%%%%%%%%%%%%%%%%%%%%%
\question \textbf{Twist on the display.} A capture of the `9' key has a row tone of
amplitude 0.20 and a column tone of amplitude 0.12.
\begin{parts}

	\begin{minipage}[t]{0.58\linewidth}
		\part In \texttt{L13-dtmf-rx.grc}, each stock Goertzel block feeds a Complex to Mag,
		which shows half a tone's amplitude. What heights do the 852 Hz and 1477 Hz traces
		reach?
		\begin{solution}[0.7in]
		The 852 Hz trace reaches 0.10 and the 1477 Hz trace 0.06.
		\end{solution}
	\end{minipage}\hfill%
	\begin{minipage}[t]{0.37\linewidth}
		\raggedleft \vspace{12pt}
		852: \ansbox[1.0in]{0.10}

		\vspace{4pt}
		1477: \ansbox[1.0in]{0.06}
	\end{minipage}

	\part Which tone arrives weaker? Which slider in \texttt{L13-dtmf-link.grc} makes the
	same kind of signal, and what happens to the Digits label as you push that slider
	further the same way?
		\begin{solution}[1.0in]
		The column tone, 1477 Hz, is weaker. The Twist slider sets it; negative twist makes the
		column tone weaker. Digits keep decoding for a while, then stop at about $-10$ dB, the
		edge of test 3.
		\end{solution}

	\part Why does the column tone often arrive weaker over a telephone line? Answer in
	one sentence.
		\begin{solution}[0.6in]
		A telephone line loses more at high frequencies, and the column tones are the higher
		group.
		\end{solution}

\end{parts}

\newpage
%%%%%%%%%%%%%%%%%%%%%%%%%%%%%%%%%%%%%%%%%
\question \textbf{The decision rules.} The values below are one block's shares of its
energy in each bin, as the L13 decoder sees them. Its tests: (1) the best row and column
bins together hold at least 0.35; (2) each beats every other bin in its group by 6 dB,
about four times the share; (3) the column share is between a tenth and four times the
row share; (4) three blocks in a row agree.

\begin{center}
\renewcommand{\arraystretch}{1.2}\setlength{\tabcolsep}{4pt}\small
\begin{tabular}{| l | c | c | c | c | c | c | c | c |}
\hline
Bin (Hz) & 697 & 770 & 852 & 941 & 1209 & 1336 & 1477 & 1633 \\\hline
Block A & 0.02 & 0.41 & 0.03 & 0.01 & 0.01 & 0.02 & 0.38 & 0.02 \\\hline
Block B & 0.12 & 0.10 & 0.02 & 0.01 & 0.03 & 0.30 & 0.02 & 0.01 \\\hline
\end{tabular}
\end{center}

\begin{parts}

	\part Apply tests 1 to 3 to block A. Which digit does block A hold?
		\begin{solution}[1.1in]
		The best row bin is 770 Hz and the best column bin is 1477 Hz. Test 1:
		$0.41 + 0.38 = 0.79$, above 0.35. Test 2: 0.41 is far more than four times the next
		row bin, 0.03, and 0.38 far more than four times 0.02. Test 3: 0.38 is close to 0.41,
		well inside the limits. Block A holds `6'.
		\end{solution}

	\part Which test does block B fail?
		\begin{solution}[0.8in]
		The share is $0.12 + 0.30 = 0.42$, enough, but the best row bin, 0.12, is barely
		larger than the next, 0.10, nowhere near four times it, so block B fails test 2 and
		holds no digit.
		\end{solution}

	\part Block A alone does not produce a digit. What else must happen before the decoder
	reports `6', and what does that requirement reject?
		\begin{solution}[0.9in]
		Three blocks in a row, each starting 102 samples after the last, must all hold `6'.
		The tone must outlast about one 25.6 ms block, so short bursts of speech or noise that happen to
		pass the tests in one block do not count as digits.
		\end{solution}

\end{parts}

\newpage
%%%%%%%%%%%%%%%%%%%%%%%%%%%%%%%%%%%%%%%%%
\question \textbf{Noise.} You run \texttt{L13-dtmf-link.grc} and lower the SNR slider from
20 dB to 0 dB.
\begin{parts}

	\part Describe what you see on the spectrum and in the Numbers label as the SNR comes
	down.
		\begin{solution}[1.0in]
		The noise floor rises across the whole 0 to 4 kHz band until the two tones barely
		stand above it. The Numbers label stays at every number correct down to about 2 dB;
		near 1 dB some numbers come out short or wrong, and at 0 dB almost none arrive.
		\end{solution}

	\part At 3 dB the tones barely show above the noise on the spectrum, yet every number
	still decodes. Explain in words why the decoder can still hear them.
		\begin{solution}[1.0in]
		With bins 39 Hz apart, nulls 39 Hz either side, each Goertzel bin sees only a small slice of the
		noise spread across the 4 kHz band, while it still sees the whole of its tone. Inside
		its bin, the tone stands well above the noise.
		\end{solution}

\end{parts}

\end{questions}
